\section*{Introduction}

Network analysis is a critical tool for capturing dependencies inherent to the development of social phenomenon. The concept of social network analysis began as early as the 1930s with Jacob Moreno's \textit{Who Will Survive?} and has developed into a rich and diverse field. In Political Science, network analysis has now been utilized to study a diverse range of topics: trade, intergovernmental organizations, sanctions, internal conflict, political behavior.\footnote{For examples, see \cite{dorussen2010}, \cite{cao2009}, \cite{cranmer2014} and for an overview of network analysis in Political Science see \cite{ward2011}.} Yet, unfortunately, a majority of studies in Political Science still rely on the standard dyadic framework. We argue that conflict evolution is a process conditioned on the relative effects of actors' behavior and is thus best conceptualized via a network approach. 

Sociologists have long established that to understand an actor's behavior you have to understand the context in which they operate as well as interpret their interactions with one partner in light of all interactions across all other partners. Conflict scholars are interested in questions that mirror these exact patterns, such as: which actor is driving the violence within a multi-actor conflict? Did a government crackdown cause anti-government coordination or chaos between armed groups? Are civilian challenges towards non-state armed actors causing an increase or decrease in violence? These questions precisely illuminate the need to apply the appropriate methodological approach to pressing questions in our field.

Our paper, however, is not merely a methodological exercise. We are also motivated by a rich, underdeveloped theoretical question: do nonviolent protests influence the evolution of violence between armed actors?  In this paper we utilize a new database of conflictual events in Mexico. We investigate whether civilian organized via nonviolent protests can influence conflict between armed actors. To do so, we generate a conflict network for each year between 2006-2010 and incorporate novel protest data which documents the number of protests against violence in municipalities across Mexico during the drug war.

